% Part: sets-functions-relations
% Chapter: relations
% Section: reflections
%
\documentclass[../../../include/open-logic-section]{subfiles}

\begin{document}

\olfileid{sfr}{rel}{ref}
\olsection{Even filosofisch nadenken}

In \olref[set]{sec} hebben we relaties gedefinieerd als bepaalde
verzamelingen. Laten we even stilstaan bij een filosofische vraag: wat
\emph{doen} we eigenlijk wanneer we zo'n definitie geven? Het is maar
de vraag of we daarmee willen zeggen dat we hebben \emph{ontdekt}
dat bepaalde dingen in werkelijkheid precies hetzelfde zijn;
bijvoorbeeld dat de ordeningsrelatie op $\Nat$ \emph{bleek} samen te
vallen met de verzameling $R=
\Setabs{\tuple{n,m}}{n, m \in \Nat\text{ en } n < m}$ die we in
\olref[set]{sec} hebben gedefinieerd. Daar zijn drie
redenen voor.

Ten eerste: in \olref[set][pai]{wienerkuratowski} kozen we
$\tuple{a, b} = \{\{a\}, \{a, b\}\}$ als definitie. Kijk nu naar de
andere definitie $\lVert a, b\rVert = \{\{b\}, \{a, b\}\} =
\tuple{b,a}$. Als $a \neq b$, dan is $\tuple{a, b} \neq \lVert
a,b\rVert$. Maar we hadden $\lVert a,b\rVert$ net zo goed in plaats
van $\tuple{a,b}$ kunnen gebruiken als definitie van een geordend paar.
Beide keuzes hadden even goed gewerkt. We hebben dus twee
verzamelingen die allebei even goed dienst kunnen doen als de
ordeningsrelatie op de natuurlijke getallen:
\begin{align*}
		R &= \Setabs{\tuple{n,m}}{n, m \in \Nat \text{ en }n < m}\\
		S &= \Setabs{\lVert n,m\rVert}{n, m \in \Nat \text{ en }n < m}.
\end{align*}
Volgens het extensionaliteitsbeginsel zijn verzamelingen gelijk als ze
dezelfde elementen hebben. Omdat $R \neq S$, kunnen ze dus niet
\emph{allebei} de ordeningsrelatie op~$\Nat$ zijn.
Maar dan zou het zomaar een keuze zijn om te zeggen dat $R$ en niet $S$
(of omgekeerd) de ordeningsrelatie \emph{is}. Dat zou nogal gênant
zijn. (Dit is een heel eenvoudig voorbeeld van een argument tegen het
idee dat je zulke wiskundige dingen volledig tot verzamelingen kunt
terugbrengen. Het argument is door Benacerraf in
\citeyear{Benacerraf1965} bekend geworden. We komen hier nog een paar
keer op terug.)

Ten tweede: stel dat \emph{elke} relatie een verzameling is. Dan moet
ook de relatie ``is element van'', $\in$, zelf één bepaalde
verzameling zijn. Dat zou de verzameling
$\Setabs{\tuple{x,y}}{x \in y}$ moeten zijn. Maar bestaat die
verzameling wel? Door de paradox van Russell kunnen we dat niet zomaar
aannemen. Sterker nog,
\oliflabeldef{cumul:::part}{de verzamelingenleer die we in
\olref[cumul][][]{part} opbouwen, zal zeggen dat deze verzameling
\emph{niet bestaat}.\footnote{We lopen even vooruit. Neem voor een
bewijs uit het ongerijmde aan dat $I =
\Setabs{\tuple{x,y}}{x \in y}$ bestaat. Dan is $\bigcup \bigcup I$ de
universele verzameling, dus de verzameling van alle verzamelingen. Dat
is in strijd met \olref[sfr][z][sep]{thm:NoUniversalSet}.}}{het mogelijk
is de verzamelingenleer nauwkeurig op te bouwen met axioma's, en in die
theorie bestaat deze verzameling inderdaad niet.}
Dus ook als we sommige relaties als verzamelingen kunnen gebruiken,
moeten we met de relatie ``is element van'' anders omgaan.

Ten derde: toen we relaties met verzamelingen ``gelijkstelden'',
spraken we af dat we $Rxy$ mochten schrijven voor $\tuple{x,y} \in R$.
Dat werkt zolang we ``$\in$'' \emph{als} een predicaat gebruiken: als
iets waarmee we zeggen dat het ene ding een element van het andere is.
Maar stel dat ``$\in$'' zelf een verzameling aanduidt. Dan is de
uitdrukking ``$\tuple{x,y} \in R$'' alleen een rij van drie losse
uitdrukkingen die elk voor één verzameling staan: ``$\tuple{x,y}$'', ``$\in$''
en ``$R$''. Met een rij namen kun je net zomin een uitspraak doen als
met deze onzin: ``het kopje pennenhouder de tafel''. Ook hier moeten we
dus anders omgaan met de relatie ``is element van'', zelfs als we
sommige andere relaties wel als verzamelingen kunnen gebruiken. (Hier
komen een eenvoudige versie van de paradox van Frege over het begrip
\emph{paard} en een bekend bezwaar van Wittgenstein tegen Russell
samen.)

Wat volgt hier nu uit? Er is niets \emph{mis} met de uitspraak dat
relaties tussen getallen verzamelingen zijn. We moeten alleen weten hoe
die uitspraak bedoeld is. We zeggen niet dat we een feit hebben ontdekt
over welke dingen in werkelijkheid één en hetzelfde zijn. We zeggen alleen dat
we in sommige situaties bepaalde relaties als bepaalde verzamelingen
kunnen en zullen \emph{behandelen}.
\end{document}